\section{引言}\label{sec:introduction}

哥德巴赫猜想是数论中最古老且最著名的未解决问题之一，其历史可追溯至18世纪。该猜想由普鲁士数学家克里斯蒂安·哥德巴赫（Christian Goldbach）于1742年在与莱昂哈德·欧拉（Leonhard Euler）的通信中首次提出。

\subsection{问题陈述}

哥德巴赫猜想通常分为两个版本：

\begin{definition}[强哥德巴赫猜想]
    每个大于2的偶数都可以表示为两个素数之和。
\end{definition}

\begin{definition}[弱哥德巴赫猜想]
    每个大于5的奇数都可以表示为三个素数之和。
\end{definition}

强哥德巴赫猜想可表述为：对于任意偶数 $n > 2$，存在素数 $p$ 和 $q$ 使得
\begin{equation}
    n = p + q
\end{equation}

弱哥德巴赫猜想可表述为：对于任意奇数 $n > 5$，存在素数 $p_1, p_2, p_3$ 使得
\begin{equation}
    n = p_1 + p_2 + p_3
\end{equation}

\subsection{研究意义}

哥德巴赫猜想的研究具有多重意义：

\begin{enumerate}[label=(\arabic*)]
    \item \textbf{理论价值}：该猜想揭示了素数分布的深层规律，与黎曼假设等核心问题密切相关。
    \item \textbf{方法论价值}：为解决该猜想而发展出的筛法、圆法等工具已成为解析数论的标准方法。
    \item \textbf{历史价值}：该问题推动了20世纪解析数论的蓬勃发展。
\end{enumerate}

\subsection{本文结构}

本文组织如下：\cref{sec:history} 回顾哥德巴赫猜想的历史发展；\cref{sec:methods} 介绍主要研究方法；\cref{sec:results} 总结重要研究成果；\cref{sec:applications} 讨论相关应用；\cref{sec:open_problems} 列举开放问题；\cref{sec:conclusion} 总结全文。
